\documentclass[12pt, a4paper]{article}
\usepackage[UTF8, zihao=-4]{ctex}
\usepackage[top=2.5cm, bottom=2.5cm, left=2.8cm, right=2.8cm]{geometry}
\usepackage{fontspec}
\setmainfont{Times New Roman}
\usepackage{setspace}
\setstretch{1.5}
\usepackage{booktabs}
\usepackage{array}
\usepackage{makecell}
\usepackage{multirow}
\usepackage{adjustbox}
\usepackage{amsmath, amssymb}
\usepackage{float}
\usepackage{caption}
\usepackage{fancyhdr}
\setlength{\headheight}{15pt}
\pagestyle{fancy}
\fancyhf{}
\fancyhead[C]{\small 实验数据记录单}
\fancyfoot[C]{\small\thepage}

\begin{document}

\begin{center}
  {\heiti\zihao{3} 星点法测量透镜像差综合实验 \\ 实验数据记录单} \\[1cm]
  \begin{tabular}{llll}
    姓\quad 名：& \underline{\hspace{4cm}} &
    学\quad 号：& \underline{\hspace{4cm}} \\[0.5cm]
    班\quad 级：& \underline{\hspace{4cm}} &
    实验日期：  & \underline{\hspace{4cm}} \\[0.5cm]
    指导教师：  & \underline{\hspace{4cm}} &
    教师签名：  & \underline{\hspace{4cm}} \\
  \end{tabular}
\end{center}
\vspace{1cm}

\subsection*{表 1 \quad 位置色差记录}
\begin{table}[H]
  \centering
  \begin{adjustbox}{max width=\textwidth}
    \begin{tabular}{cccc}
      \toprule
      次数 & 蓝光焦点 $X_1$/mm & 红光焦点 $X_2$/mm & 位置色差 $\Delta X=X_2-X_1$/mm \\
      \midrule
      1 & & & \\
      2 & & & \\
      3 & & & \\
      \midrule
      平均值 & & & \\
      \bottomrule
    \end{tabular}
  \end{adjustbox}
\end{table}
{\footnotesize 记录说明：蓝光 LED 波长 451 nm，红光 LED 波长 690 nm；测量色差时应保持同一环带光阑。}

\vspace{0.6cm}

\subsection*{表 2 \quad 倍率色差记录}
\begin{table}[H]
  \centering
  \begin{adjustbox}{max width=\textwidth}
    \begin{tabular}{ccccc}
      \toprule
      次数 & $b$/pixel & $c$/pixel & $c-b$/pixel & 倍率色差 $2.9(c-b)$/$\mu\mathrm{m}$ \\
      \midrule
      1 & & & & \\
      2 & & & & \\
      3 & & & & \\
      \midrule
      平均值 & & & & \\
      \bottomrule
    \end{tabular}
  \end{adjustbox}
\end{table}
{\footnotesize 记录说明：CMOS 单个像素大小为 2.9 $\mu\mathrm{m}$。}

\subsection*{表 3 \quad 红色光源轴向球差记录}
\begin{table}[H]
  \centering
  \begin{adjustbox}{max width=\textwidth}
    \begin{tabular}{cccc}
      \toprule
      次数 & 最小环带焦点 $X_1$/mm & 大号环带焦点 $X_2$/mm & 轴向球差 $\Delta X=X_2-X_1$/mm \\
      \midrule
      1 & & & \\
      2 & & & \\
      3 & & & \\
      \midrule
      平均值 & & & \\
      \bottomrule
    \end{tabular}
  \end{adjustbox}
\end{table}

\subsection*{表 4 \quad 红色光源垂轴球差记录}
\begin{table}[H]
  \centering
  \begin{adjustbox}{max width=\textwidth}
    \begin{tabular}{ccccc}
      \toprule
      次数 & $b$/pixel & $c$/pixel & $c-b$/pixel & 垂轴球差 $2.9(c-b)$/$\mu\mathrm{m}$ \\
      \midrule
      1 & & & & \\
      2 & & & & \\
      3 & & & & \\
      \midrule
      平均值 & & & & \\
      \bottomrule
    \end{tabular}
  \end{adjustbox}
\end{table}
{\footnotesize 记录说明：垂轴球差由 CMOS 像素坐标差换算，像素尺寸为 2.9 $\mu\mathrm{m}$。}

\subsection*{表 5 \quad 像散数据记录}
\begin{table}[H]
  \centering
  \begin{adjustbox}{max width=\textwidth}
    \begin{tabular}{cccc}
      \toprule
      次数 & $X_1$/mm & $X_2$/mm & 透镜像散 $\Delta X=X_2-X_1$/mm \\
      \midrule
      1 & & & \\
      2 & & & \\
      3 & & & \\
      \midrule
      平均值 & & & \\
      \bottomrule
    \end{tabular}
  \end{adjustbox}
\end{table}

\noindent\textbf{慧差观察记录：}

\vspace{2cm}

\noindent\textbf{其它现象与误差来源记录：}

\vspace{2cm}

\end{document}
